\setcounter{subsection}{0}

\section{CONCLUSIONES}
\sloppy

\subsection{Análisis de impacto del proyecto}
Para un equipo que integra pagos en Colombia, el impacto inmediato es práctico. Un paquete de código abierto, con licencia Apache 2.0, cubre tarjeta y PSE en cuatro pasarelas, y la API de Simulación permite probar rechazos, pagos pendientes y webhooks sin credenciales reales. La documentación de datos de prueba sirve también a quien no use el framework, porque reúne lo que acepta y lo que rechaza cada sandbox. La conclusión de OE-06 y el capítulo VII responden cuánto reduce el framework el esfuerzo de integración.

\subsubsection{Impacto en la ingeniería de sistemas}
Para la Ingeniería de Sistemas, el proyecto muestra una forma de aplicar Design Science Research a un framework de software con la evidencia versionada junto al código. Los requisitos están trazados hasta sus pruebas, las decisiones quedaron registradas y los documentos compilan en la misma integración continua.

\subsubsection{Impacto en contextos globales, económicos, ambientales y sociales}
El proyecto no cambió las condiciones del ecosistema de pagos colombiano. No interviene en los costos de procesamiento de las pasarelas ni resuelve problemas de acceso financiero. Su aporte es técnico y se limita a la fragmentación del proceso de integración. A corto plazo, el beneficio recae en desarrolladores individuales y equipos pequeños, que encuentran en un solo punto de entrada información y lógica que estaba dispersa. A mediano plazo, si el framework se adopta fuera de este proyecto, puede llegar a la comunidad que atiende el comercio electrónico colombiano y a los programas académicos que enseñen integración de sistemas, para los que el repositorio sirve como un caso concreto.

A largo plazo, el impacto depende de condiciones que el proyecto no controla: que alguien mantenga el framework, que las APIs de las pasarelas no cambien y que la comunidad lo adopte. Si se cumplen, emprendimientos y pequeñas empresas podrían incorporar varias opciones de cobro digital sin duplicar esfuerzo, lo que reduciría de forma marginal las barreras de entrada al comercio electrónico formal. Ese efecto social es una expectativa condicionada y no un resultado medido. No se midieron efectos económicos ni ambientales, y por eso no se afirma ninguno.

\subsection{Conclusiones y trabajo futuro}
\subsubsection{Cumplimiento de los objetivos y contribuciones}
Este apartado responde a cada objetivo específico con la evidencia del repositorio y con la validación del capítulo VII. El estado de cada uno es el que registran el seguimiento de los objetivos (Artefacto~\ref{anx:objetivos}) y la matriz de trazabilidad (Artefacto~\ref{anx:trazabilidad}), que da por implementados los diecinueve requisitos funcionales.

\paragraph{OE-01 (análisis de las pasarelas).} Se cumplió. El análisis está en la caracterización técnica de las cuatro pasarelas (Artefacto~\ref{anx:producto}), en el glosario del dominio (Artefacto~\ref{anx:arquitectura}) y en la documentación de datos de prueba (Artefacto~\ref{anx:datos}), y el capítulo III lo resume. Mostró que las cuatro pasarelas difieren justo en lo que más pesa al integrarlas. Wompi firma con un \textit{hash} SHA-256 y las otras tres con HMAC-SHA256, y la firma de Mercado Pago no cubre el cuerpo del webhook. Kushki rechaza una tarjeta con un \texttt{400} en lugar de una transacción rechazada, y Rapyd firma cada petición y cobra la tarjeta en su página. También mostró un límite: los sandboxes de Wompi, Mercado Pago y Kushki no permiten probar varios flujos de punta a punta, y esos flujos se prueban con la API de Simulación. La validación confirmó el análisis con las cuatro pasarelas reales, porque las 25 pruebas de contrato pasaron en sus sandboxes (Tabla~\ref{tab:suites-memoria}). DEF-01, uno de los dos defectos de la corrida, corrigió en la documentación el estado nativo de Mercado Pago, que es \texttt{processed} y no \texttt{approved} (capítulo VII, apartado~\ref{sec:vii-demostracion}).

\paragraph{OE-02 (diseño de la arquitectura).} Se cumplió. El SDK aplica la arquitectura hexagonal: el dominio no depende de ningún adaptador, y cada pasarela es un adaptador de \texttt{PaymentGatewayPort}. La fachada \texttt{KitPagos}, la fábrica \texttt{GatewayFactory} y las estrategias de normalización vienen del catálogo GoF (capítulo VI; Equipo del Proyecto, 2026a). Las decisiones quedaron en el registro de decisiones de arquitectura (Artefacto~\ref{anx:registro}). La integración continua verifica en cada cambio los umbrales de WMC, CBO, RFC y complejidad máxima del SDK y de la API de Simulación, aunque no todas las clases caben en ellos. El SDK declara tres excepciones: el CBO de \texttt{Transaction}, el WMC de \texttt{Amount} y el CBO que el tipo \texttt{PseBank} suma a cinco clases. La API de Simulación declara una, el CBO de \texttt{KitPagosProvider}. Cada excepción tiene su justificación escrita en el registro, como exige la metodología del proyecto (Artefacto~\ref{anx:metodologia}). Con ellas, las 35 clases del SDK y las 28 del simulador quedan dentro de los umbrales (capítulo VII, apartado~\ref{sec:vii-pruebas}).

\paragraph{OE-03 (prototipo funcional del framework).} Se cumplió. El SDK está publicado en npm, con la 0.6.0 en \texttt{latest} y la 0.7.0-rc.2 en \texttt{next} (Kit Pagos Colombia, 2026). Implementa la interfaz unificada para tarjeta y PSE en las cuatro pasarelas y verifica la firma de los webhooks de cada una en tiempo constante y con ventana de tolerancia. También reintenta las operaciones idempotentes y oculta las credenciales en sus errores (RF-01 a RF-08 y RF-16 a RF-19). La ejecución del plan de pruebas lo verificó sobre la 0.7.0-rc.2: la suite del SDK pasó sus 58 suites y 1088 pruebas, con 97,95\,\% de sentencias cubiertas (Tabla~\ref{tab:suites-memoria}). CP-01, CP-05 y CP-07 mostraron el mismo pago por las cuatro pasarelas, la tokenización en el navegador y los reintentos limitados a las lecturas (Tabla~\ref{tab:casos-memoria}). La interfaz unificada tiene límites que el SDK declara en lugar de esconder:
\begin{itemize}
    \item con Rapyd, el cobro con tarjeta es una redirección, porque Rapyd no tokeniza en el navegador;
    \item en Mercado Pago, el evento del webhook llega como \texttt{PENDING}, y el comercio lo confirma con \texttt{getPaymentStatus()};
    \item la 0.6.0 entregaba el rechazo de tarjeta de Kushki como un error técnico (\texttt{INVALID\_REQUEST}), y la 0.7.0-rc.2 lo entrega como una transacción rechazada.
\end{itemize}

\paragraph{OE-04 (API de Simulación).} Se cumplió. Bajo \texttt{/v1/sim}, la API imita las rutas y los formatos de las cuatro pasarelas. Elige el escenario por cabecera, dato de prueba, monto reservado o marca de credencial, y emite los webhooks firmados cuando un pago cambia de estado. Bajo \texttt{/v1/api} ofrece por HTTP las operaciones del framework, que el prototipo C usa desde el navegador. Se ejecuta en local y en Render, en el plan gratuito. En la validación, su suite pasó 54 suites y 1301 pruebas, y CP-06 recorrió \texttt{/v1/api} en local y con una prueba de humo en Render. La carga terminó sus 30 corridas sin errores (capítulo VII, apartado~\ref{sec:vii-pruebas}). Quedan dos límites declarados: las combinaciones sin evidencia de cómo responde la pasarela real devuelven \texttt{501}, y siete de sus 43 operaciones no tienen esquema OpenAPI.

\paragraph{OE-05 (datos de prueba).} Se cumplió. La documentación de datos de prueba (Artefacto~\ref{anx:datos}) reúne, por pasarela, las credenciales de sandbox, las tarjetas, los documentos, los bancos y los flujos medidos. Cada dato declara su nivel de evidencia: medido en la API real con fecha, medido hasta cierto punto del flujo o tomado de la documentación sin comprobar. La API de Simulación reproduce esos datos, salvo los de Rapyd y los del PSE de Mercado Pago, donde decide la página del banco o de la pasarela. Las tarjetas que usan los casos de prueba salen de esa documentación. CP-02 eligió el banco de prueba que aprueba en Wompi, y CP-03 reprodujo los rechazos de Kushki y Mercado Pago (capítulo VII, apartado~\ref{sec:vii-demostracion}).

\paragraph{OE-06 (evaluación con los prototipos).} Se cumplió. El informe comparativo de la Fase 5 (Artefacto~\ref{anx:informe}) compara el prototipo B, que integra las cuatro pasarelas con el SDK, con el D, que las integra a mano. La comparación usa las seis variables del plan, medidas entre el 22 y el 28 de octubre de 2026 (Tabla~\ref{tab:seis-variables-memoria}). En el módulo de pagos, B tiene 310 líneas de código de integración frente a 1\,308 de D, un 76,30\,\% menos, y sumar las tres pasarelas nuevas le costó 67 líneas agregadas frente a 1\,042. Su WMC total, que suma la complejidad de todas sus funciones, baja de 375 a 66, y la complejidad de su función más compleja, de 19 a 10. B no expone ningún concepto nativo de pasarela, frente a 10 identificadores en D, alcanza las mismas cuatro pasarelas y cubre el 100\,\% de sus sentencias, frente al 90,03\,\% de D (capítulo VII, apartado~\ref{sec:vii-evaluacion}).

Dos resultados no favorecen al framework, y el informe y el capítulo VII los reportan tal como salieron (apartado~\ref{sec:vii-discusion}). El CBO promedio de B es mayor, 7,0 frente a 5,67, porque su módulo concentra más tipos externos por archivo. Además, 9 de sus 67 líneas no son configuración, porque Kushki y Rapyd obligaron a ajustar el módulo común. La medición cuenta el código del comercio y no el costo de construir el SDK, y es un estudio de caso de dos prototipos escritos por el mismo equipo.

Además del framework, el trabajo deja dos prácticas que se pueden reutilizar. La primera es un registro de decisiones que separa, para cada pasarela, lo medido en el sandbox de lo que dice su documentación. La segunda son las métricas CK verificadas en la integración continua como parte de la DoD, con cada excepción justificada por escrito.

\subsubsection{Trabajo futuro}
El trabajo futuro sale de los límites que el proyecto declaró:
\begin{itemize}
    \item \textbf{Publicar la 0.7.0 como versión estable.} Hoy está como 0.7.0-rc.2 en la etiqueta \texttt{next}, y la 0.6.0 de \texttt{latest} todavía entrega el rechazo de tarjeta de Kushki como un error técnico.
    \item \textbf{Completar P3 con Rapyd} si Rapyd llega a ofrecer tokenización en el navegador. Mientras tanto, el cobro con tarjeta de Rapyd sigue en su página alojada.
    \item \textbf{Dar esquema a las siete operaciones de la API de Simulación que no lo tienen}.
    \item \textbf{Ampliar los métodos de pago} a los que el alcance dejó fuera: efectivo, Nequi, Daviplata, suscripciones y reembolsos.
    \item \textbf{Probar de punta a punta los flujos que los sandboxes no permiten}, cuando las pasarelas lo habiliten, y reemplazar con evidencia medida las combinaciones que hoy responden \texttt{501}.
    \item \textbf{Repetir la evaluación con autores distintos para B y D y con más de dos prototipos.} Un solo autor escribió las pasarelas nuevas en los dos, y dos prototipos no son una muestra (Artefacto~\ref{anx:informe}, secciones 5.3 y 5.7).
\end{itemize}
